\refstepcounter{section}
\section*{Lecture 03: Gates}
\addcontentsline{toc}{section}{Lecture 03: Gates}
Having seen the basic building blocks for quantum computation, let us see how the computation can be actually carried out. A quantum computer is expected to necessarily satisfy the following conditions put forth by \textit{David P. DiVincenzo} in 1996:
\begin{itemize}
    \item A scalable system with well-characterised qubits, that is, the system should be able to accomodate enough number of qubits necessary for computation. 
    \item Initialising the qubits to some simple \textit{fiducial state}, so that computations can be carried forward from there.
    \item Long relevant \textit{coherence time}, that is, the system does not undergo decoherence until computation is finished. 
    \item A universal set of quantum gates, which will carry out the computation by evolving the states. These gates need to be \textit{unitary} as required by probability conservation.
    \item A specific measurement capability, that is, there should be provision to read out the results after computation is done. 
\end{itemize}
Classical computations are carried out by combination of different logic gates which are based on Boolean algebra. \\[0.2cm]
Boolean Algebra deals with logical operations of binary variables. A Boolean function $f$ is generally defined as $f: \{0,1\}^k \rightarrow \{0,1\} $, that is, it takes input from $k$ cartesian products of the set $\{0,1\}$ and returns a single output which can be 0 or 1. The basic boolean functions are: 
\begin{itemize}
    \item \textbf{NOT Gate}\\[0.3cm]
    It is a unary operator which takes a single input and returns the negation of the input. It is represented by the symbol $\overline{\mathrm{A}}$. If input is 1, it return 0 and if input is 0, it returns 1.
    \begin{figure}[H]
        \centering
        \begin{circuitikz}
            \draw (0,0) node[ieeestd not port, fill=cyan!40!white] (mynot) {};
            \draw (mynot.in) -- ++(-1,0) node[left] {A};
            \draw (mynot.out) -- ++(1,0) node[right] {Y};
        \end{circuitikz}
        \caption{Representation of NOT Gate}
    \end{figure}
    \item \textbf{AND Gate}\\[0.3cm]
    It is a binary operator which takes two inputs and returns the logical AND of the inputs. It is represented by the symbol $\mathrm{A}\cdot \mathrm{B}$. It returns 1 only if both inputs are 1, otherwise it returns 0.
    \begin{figure}[H]
        \centering
        \begin{circuitikz}
            \draw (0,0) node[ieeestd and port, fill=cyan!40!white] (myand) {};
            \draw (myand.in 1) -- ++(-1,0) node[left] {A};
            \draw (myand.in 2) -- ++(-1,0) node[left] {B};
            \draw (myand.out) -- ++(1,0) node[right] {Y};
        \end{circuitikz}
        \caption{Representation of AND Gate}
    \end{figure}
    \item  \textbf{OR Gate}\\[0.3cm]
    It is a binary operator which takes two inputs and returns the logical OR of the inputs. It is represented by the symbol $\mathrm{A} + \mathrm{B}$. It returns 1 if any of the inputs is 1, otherwise it returns 0.
    \begin{figure}[H]
        \centering
        \begin{circuitikz}
            \draw (0,0) node[ieeestd or port, fill=cyan!40!white] (myor) {};
            \draw (myor.in 1) -- ++(-1,0) node[left] {A};
            \draw (myor.in 2) -- ++(-1,0) node[left] {B};
            \draw (myor.out) -- ++(1,0) node[right] {Y};
        \end{circuitikz}
        \caption{Representation of OR Gate}
    \end{figure}
    \item \textbf{XOR Gate}\\[0.3cm]
    It is a binary operator which takes two inputs and returns the logical XOR of the inputs. It is represented by the symbol $\mathrm{A} \oplus \mathrm{B}$. It can be showed that the expression for XOR is equivalent to $\mathrm{\overline{A}B+A\overline{B}}$. It returns 1 if the inputs are different, otherwise it returns 0.
    \begin{figure}[H]
        \centering
        \begin{circuitikz}
            \draw (0,0) node[ieeestd xor port, fill=cyan!40!white] (myxor) {};
            \draw (myxor.in 1) -- ++(-1,0) node[left] {A};
            \draw (myxor.in 2) -- ++(-1,0) node[left] {B};
            \draw (myxor.out) -- ++(1,0) node[right] {Y};
        \end{circuitikz}
        \caption{Representation of XOR Gate}
    \end{figure}
    \item \textbf{NAND Gate}\\[0.3cm]
    It is a binary operator which takes two inputs and returns the logical NAND of the inputs. It is represented by the symbol $\overline{\mathrm{A} \mathrm{B}}$. Thus, it is a composition of AND and NOT operations. From de Morgan's law, it can be shown that NAND is equivalent to $\overline{\mathrm{A}} + \overline{\mathrm{B}}$. 
    \begin{figure}[H]
        \centering
        \begin{circuitikz}
            \draw (0,0) node[ieeestd nand port, fill=cyan!40!white] (mynand) {};
            \draw (mynand.in 1) -- ++(-1,0) node[left] {A};
            \draw (mynand.in 2) -- ++(-1,0) node[left] {B};
            \draw (mynand.out) -- ++(1,0) node[right] {Y};
        \end{circuitikz}
        \caption{Representation of NAND Gate}
    \end{figure}
    \item \textbf{NOR Gate}
    It is a binary operator which takes two inputs and returns the logical NOR of the inputs. It is represented by the symbol $\overline{\mathrm{A} + \mathrm{B}}$. Thus, it is a composition of OR and NOT operations. From de Morgan's law, it can be shown that NOR is equivalent to $\overline{\mathrm{A}} \cdot \overline{\mathrm{B}}$.
    \begin{figure}[H]
        \centering
        \begin{circuitikz}
            \draw (0,0) node[ieeestd nor port, fill=cyan!40!white] (mynor) {};
            \draw (mynor.in 1) -- ++(-1,0) node[left] {A};
            \draw (mynor.in 2) -- ++(-1,0) node[left] {B};
            \draw (mynor.out) -- ++(1,0) node[right] {Y};
        \end{circuitikz}
        \caption{Representation of NOR Gate}
    \end{figure}
\end{itemize}
Note that the except for the NOT gate, the number of output in each gate is always lesser than the number of input. This reduction of the computational \textit{phase space complexity} renders these classical gates irreversible, since we cannot reconstruct the input given a particular output. This is problematic since we require quantum logic gates to be reversible due to unitarity condition, so a different approach is to be taken.\\[0.2cm]
\textbf{NOTE:} Not all classical gates are irreversible. For example, the Toffoli gate is a universal reversible gate. It is three-input three-output controlled gate which means that the third input (`target') is controlled by the first two inputs (`control'), both of which do not change. That is, the target bit will be inverted if the first and second bits are both $1$. Given inputs $a,b,c$ then the output is accordingly given as $y=c\oplus (a\cdot b)$.

\begin{table}[H]
    \centering
    \caption{Truth Table for the Toffoli gate}
    \vspace{0.2cm}
    \begin{tabular}{|c|c|c||c|}
    \hline
    \textbf{A} & \textbf{B} & \textbf{C} & \textbf{Output (a,b,y)}  \\
    \hline
    0 & 0 & 0  & 0,0,0 \\
    0 & 0 & 1  & 0,0,1 \\
    0 & 1 & 0 & 0,1,0 \\
    0 & 1 & 1  & 0,1,1 \\
    1 & 0 & 0  & 1,0,0 \\
    1 & 0 & 1  & 1,0,1 \\
    1 & 1 & 0 & 1,1,1 \\
    1 & 1 & 1  & 1,1,0 \\
    \hline
    \end{tabular}
    \end{table}
    \noindent
    There are many different quantum gates, as we will see\footnote{The title page itself shows two such gates.}. One of the most useful gate is the single-qubit \textit{Hadamard gate} represented as,
     \begin{figure}[H]
            \centering
            \begin{tikzpicture}
\node[scale=1.2]{
            \begin{quantikz}
    {\ket{\psi}\ } & \gate[style={fill=cyan!20,rounded
corners}, label style= black]{H} &\ \ket{\phi}
    \end{quantikz}};\end{tikzpicture}
        \end{figure} 
        \noindent
The Hadamard gate is a unitary gate and is defined by its action on the basis elements:
$$\ket{0}\xmapsto{H}\ket{+} = \frac{1}{\sqrt{2}}(\ket{0}+\ket{1})\qquad \ket{0}\xmapsto{H}\ket{-}= \frac{1}{\sqrt{2}}(\ket{0}-\ket{1})$$
Note that $H\ket{+} = \frac{1}{\sqrt{2}}(H\ket{0}+H\ket{1}) = \frac{1}{\sqrt{2}}(\ket{+}+\ket{-}) = \ket{0}$. Similarly, $H\ket{-}=\ket{1}$. Thus, $H^2=\iden$ which implies that $H=H^{-1}$. This is also evident from the Hadamard matrix written using the computation basis\footnote{To obtain the matrix element in a particular basis, we have $H_{ij} =\braket{i|H|j}$. So let's say $H_{00}$ for the Hadamard gate will be, $$\braket{0|H|0} = \braket{0|+} = \frac{1}{\sqrt{2}}(\braket{0|0}+\braket{0|1}) = \frac{1}{\sqrt{2}}$$}, 
\begin{align*}
    H\equiv \frac{1}{\sqrt{2}}\mqty(1& 1 \\ 1 & -1)
\end{align*}
    \subsection{Deutsch's Algorithm}
    Let us demonstrate a quantum algorithm which demonstrates `quantum parallelism'. Parallelism is meant to convey the sense that multiple computations can be simultaneously performed which is not possible in classical computation. 